\section{The moving-region interpretation of $A-T$}
\label{sec:moving-region-AT}
%======================================================================

This section translates the algebraic identity $A-T$ into a geometric partition.  It is the bridge between the original Mertens sum and the later kernel formulation.

Using the signed measure \cref{eq:signed-measure}, define
\begin{equation}
 \mathcal A(R):=\nu(\mathscr A_R),
 \qquad
 \mathcal U(R):=\nu(\mathscr U_R).
 \label{eq:AU-measures}
\end{equation}
At $R=R_n$,
\begin{equation}
 \mathcal A(R_n)=A_n,
 \qquad
 \mathcal U(R_n)=-T_n.
 \label{eq:AU-AT}
\end{equation}
Indeed, the transport points satisfy $c<R_n<q$ and carry signed point mass
\[
 \mu(cq)=-\mu(c).
\]
Therefore
\begin{equation}
 \boxed{
 S_n
 =\nu(\mathscr P_{R_n})
 =\mathcal A(R_n)+\mathcal U(R_n)
 =A_n-T_n.
 }
 \label{eq:geometric-AT}
\end{equation}

Equivalently, define projectors
\begin{align}
 P_R(w)&:=\1_{\{\Repart w<R^2\}},
 \label{eq:P-projector}\\
 Q_R(w)&:=\1_{\{|w|+\Impart w\le R^2\}}.
 \label{eq:Q-projector}
\end{align}
Then
\begin{align}
 \mathcal A(R)&=\int P_RQ_R\,d\nu,
 \label{eq:A-projector}\\
 \mathcal U(R)&=\int P_R(1-Q_R)\,d\nu,
 \label{eq:U-projector}\\
 S(R)&=\int P_R\,d\nu.
 \label{eq:S-projector}
\end{align}
The identity is the pointwise partition
\begin{equation}
 P_R=P_RQ_R+P_R(1-Q_R).
 \label{eq:projector-partition}
\end{equation}

%======================================================================
